% ID: SEL-D0059-OB-E01
% ESTADO: validada. S059.
\textbf{a)} Se toma la zona con $x\ge0$, por encima de $200-2x$ y $(x-100)/3$ y por debajo de $300-x/2$. En $x=0$ quedan los extremos $(0,200)$ y $(0,300)$. Las cotas inferiores se cruzan en $(100,0)$. La cota inferior ascendente corta a la superior en $(400,100)$. El otro cruce con la superior tiene $x<0$, por lo que queda fuera.
\begin{center}
\setlength{\unitlength}{0.9mm}
\begin{picture}(100,65)(-17,-8)
\qbezier(0,0)(35,0)(70,0)
\qbezier(0,0)(0,24)(0,48)
\put(71,0){$x$}\put(2,47){$y$}
\qbezier(0,30)(0,37.5)(0,45)
\qbezier(0,45)(30,30)(60,15)
\qbezier(60,15)(37.5,7.5)(15,0)
\qbezier(15,0)(7.5,15)(0,30)
\put(-17,30){\scriptsize $(0,200)$}
\put(-17,43){\scriptsize $(0,300)$}
\put(12,-6){\scriptsize $(100,0)$}
\put(61,15){\scriptsize $(400,100)$}
\end{picture}
\end{center}
\textbf{b)} Esta segunda región es distinta. Evaluamos la función lineal en sus vértices:
\[
\begin{array}{c|rrrrr}
&A&B&C&D&E\\\hline
F&31&51&71&71&46
\end{array}
\]
El mínimo es $\boxed{31}$ sólo en $A$. El máximo es $\boxed{71}$ en todo el segmento $CD$, pues las combinaciones de $C$ y $D$ conservan el valor y cualquier combinación que incluya otro vértice con peso positivo lo reduce.
